% runtime/part_intro-DE.tex

\chapter*{Laufzeit}
\phantomsection
\addcontentsline{toc}{chapter}{Laufzeit}

Dieser Teil legt dar, wie genehmigte ingenieurtechnische Semantik operativ
bewertet wird, ohne die begriffliche Verantwortung zu übertragen.

Grundlagen und Zustand haben die persistente ingenieurtechnische Information
und die kanonischen Zustandsebenen begründet. Die Laufzeit verwendet diese
stabilisierten Begriffe und bewertet sie in operativen Kontexten. Sie
definiert weder Alignment-Semantik, pose2/pose3, metrische Realisierung,
physische Realisierung noch Darstellung neu.

Der Laufzeitteil folgt daher der kanonischen Abfolge
\[
\begin{aligned}
&\text{persistente ingenieurtechnische Information}
\rightarrow
\text{Laufzeitbewertung}
\\
&\rightarrow
\text{Laufzeitdienste}
\\
&\rightarrow
\text{abgeleiteter ingenieurtechnischer Zustand}
\\
&\rightarrow
\text{Darstellung}
\\
&\rightarrow
\substack{\text{anwendungsspezifische}\\\text{Interpretation}}.
\end{aligned}
\]

In dieser Abfolge sind Laufzeitdienste operatorbasierte Bewertungssysteme.
Sie rekonstruieren und interpretieren Größen aus bestehenden Zustands- und
Realisierungsebenen, statt neue Engineering Objects einzuführen.

\section*{Bezug zur AIM-Roadmap}

In der AIM-Roadmap entspricht dieser Teil dem Übergang von
Zustandsraummodellen zu operativer Bewertung und Interpretation.

\begin{figure}[ht]
\centering
% figures/runtime/icon_runtime_operator_chain-DE.tex
\begin{tikzpicture}[
  node distance=1.25cm,
  box/.style={draw, rounded corners, minimum width=4.6cm, minimum height=0.9cm, align=center},
  op/.style={draw, rounded corners, minimum width=4.6cm, minimum height=0.9cm, align=center},
  derived/.style={draw, rounded corners, dashed, minimum width=4.6cm, minimum height=0.9cm, align=center},
  arrow/.style={->, thick}
]
\node[box] (info) {Ingenieurtechnische Information};
\node[op, below=of info] (ops) {Operatoren};
\node[derived, below=of ops] (runtime) {Abgeleiteter Laufzeitzustand};
\draw[arrow] (info) -- (ops);
\draw[arrow] (ops) -- (runtime);
\end{tikzpicture}

\caption{Laufzeitsymbol: operatorvermittelte Ableitung von ingenieurtechnischer Information zu Laufzeitzustand.}
\label{fig:icon-runtime-operator-chain}
\end{figure}

\noindent
Die Laufzeitkapitel spezialisieren diese Operatorkette nun für konkrete
Dienste und Projektionen.

Die Leitfrage der Abbildung lautet: Welche Operatorkette muss ausdrücklich
erhalten bleiben, damit die Laufzeit reproduzierbar, interpretierbar und
architektonisch konsistent bleibt?

\begin{principle}[Prinzip der Laufzeitebene]
\label{princ:runtime-layer}
In AIM werden operative Größen durch Laufzeitoperatoren erzeugt, die auf
bestehenden Zustands- und Realisierungsebenen wirken, und nicht redundant als
unabhängige Wahrheit gespeichert.
\end{principle}

\begin{summary}
Die Laufzeitebene ist die operatorbasierte Bewertungsebene von AIM.

Sie verwendet persistente ingenieurtechnische Information, erzeugt
abgeleiteten Laufzeitzustand und ermöglicht nachgelagerte Darstellung und
anwendungsspezifische Interpretation, ohne vorgelagerte Semantik neu zu
definieren.
\end{summary}

Diese Laufzeitarchitektur bereitet den Teil Realität vor, in dem der bewertete
Zustand mit physischer Realisierung, Messung und Vergleich in Beziehung gesetzt
wird.
